\documentclass[10pt]{article}

\usepackage[margin=0.75in]{geometry}
\usepackage{amsmath}
\usepackage{times}
\usepackage{titlesec}
\usepackage{enumitem}
\usepackage{xcolor}
\usepackage{hyperref}
\hypersetup{colorlinks=true, linkcolor=blue, urlcolor=blue}

\titlespacing*{\section}{0pt}{8pt}{4pt}
\titlespacing*{\subsection}{0pt}{6pt}{2pt}
\setlist[itemize]{topsep=2pt, itemsep=1pt, parsep=0pt}

\pagenumbering{gobble}

\begin{document}

\begin{center}
{\LARGE \textbf{Wingman: A Volatility-Surface Trading Agent for SPY Options}}\\[4pt]
% TODO: confirm this is the account ID actually being submitted before
% the write-up is finalized -- flag if the parallel/backup account ends
% up being used instead.
Paper account: \texttt{PA3ICB1OURR1} \quad|\quad
% TODO: replace with the actual repository URL.
\href{https://github.com/}{github.com/[repo]}
\end{center}

\vspace{2pt}
\hrule
\vspace{6pt}

\section*{Overview}
Wingman is an autonomous SPY options-trading pipeline built on Alpaca's paper-trading
API. Each cycle it fetches a live SPY option-chain snapshot, fits a two-component
lognormal mixture model (Brigo--Mercurio) to the market-implied volatility smile,
compares the model's price at each strike against the market's, and proposes a
defined-risk structure -- a long straddle where the market is cheap versus the model,
or a short call vertical where it is rich -- subject to a chain of materiality and
liquidity gates before any order is built.
% TODO (optional, space permitting): if a final week P&L or trade-count figure is
% available after tonight's close, a one-clause addition here ("...; across N live
% cycles this week the pipeline proposed M structures totalling $X unrealized P&L")
% would fit without pushing the write-up past one page. Do not add a number until
% confirmed against the account snapshot at Thursday EOD.

\section*{AI / Decision Logic}
The core signal is a residual, not a raw IV level: implied vols are recovered from
live bid/ask mids via Black--Scholes inversion, fit in price space to the mixture
model via nonlinear least squares (Levenberg--Marquardt), and the fitted model's
price at each strike is compared against the market's call-equivalent mid. A
residual that clears every gate below becomes the highest-conviction trade
candidate for that cycle.

A planned second layer asks an LLM (Gemini) to classify the current market regime
and gate proposed trade sizing (full / half / stand-down) before execution.
% TODO: this sentence must be corrected or removed if the regime gate is still not
% implemented by submission time -- as of the last working session,
% engine/regime_gate.py raises NotImplementedError and is not yet wired into the
% live cycle. Do not let this read as "implemented" unless it demonstrably is.

\section*{Risk Gates}
Every proposal must clear all of the following before an order is built:
\begin{itemize}
  \item \textbf{Degenerate-fit guard}: both fitted volatilities must sit inside an
    economically plausible SPY vol band; a fit against its numerical bounds is
    rejected outright rather than trusted as signal.
  \item \textbf{Liquidity filter}: strikes are only considered if both legs clear a
    minimum open-interest threshold and carry a genuine two-sided quote.
  \item \textbf{Tradable moneyness band}: candidates are restricted to strikes close
    to spot, since the unshifted symmetric mixture cannot reproduce SPY's actual
    skew and would otherwise misprice the wings by construction, not by real
    mispricing.
  \item \textbf{Materiality gates}: a residual must exceed both a multiple of the
    strike's own quote half-spread and a multiple of the fit's own RMSE noise
    floor -- otherwise it is indistinguishable from quote noise or estimation
    error.
  \item \textbf{Execution-spread gate}: applied again on the \emph{actual} legs to
    be traded (which can be materially wider than the reference quote used for the
    residual), so a wide real spread cannot eat a signal that looked clean on
    paper.
  \item \textbf{Structural risk control}: every short leg is submitted only inside
    a multi-leg (mleg) order alongside its covering long leg, so a naked short
    position can never exist even transiently.
  \item \textbf{Dry-run default}: order submission defaults to a safety dry run;
    live submission requires an explicit, non-default environment flag.
\end{itemize}
% TODO: add a position-limit / max-open-structures gate here once implemented in
% engine/regime_gate.py -- informed by the actual intraday position-accumulation
% count observed in Monday's log, not an assumed number.

\section*{Alpaca Infrastructure}
Market data (spot trades, option snapshots with bid/ask/IV/greeks) and reference
data (option contracts, open interest, the market clock) are pulled via the
\texttt{alpaca-py} SDK against Alpaca's paper endpoints. Order submission runs
through the Alpaca CLI: a multi-leg order payload is written to a temp file and
piped into \texttt{alpaca api POST /v2/orders}, the confirmed-reliable path for
mleg orders on this platform. Account equity and every open position are
snapshotted once per cycle for a mark-to-market equity time series alongside the
decision log.

\vspace{4pt}
\textit{% TODO: once tonight's EOD close is available, this closing line can carry
% one concrete figure (e.g. final week P\&L, total structures opened, or final
% account equity vs. starting equity) -- confirm the number against the account
% snapshot before inserting; do not estimate.
Wingman is an in-progress build: the mixture-fit and gating pipeline runs live and
has produced real (paper) decisions this week, while the LLM-based regime gate is
still under active development.}

\end{document}
